\documentclass[a4paper,10pt]{article}
\usepackage[francais]{babel}
\usepackage{url}

\selectlanguage{francais}
../letter-ensiie.sty

\selectlanguage{francais}
\enCouleur

\begin{document}


\name{St\'ephane \textsc{Genaud}}
\fonction{Professor at ENSIIE}
\telephoneillkirch{+33 (0)3688 54542}
\email{genaud@ensiie.fr}

\begin{letter}[Recommendation de Julien Casarin]{à}%
{Ecole Doctorale MSII}

\vspace{-1cm}
Chers Collègues,

Je connais  Julien Casarin depuis  sa première  année de scolarité  dans l'école
d'ingénieur  en  informatique   ENSIIE,  dont  je  dirige   l'antenne  locale  à
Strasbourg, et dans  laquelle j'interviens également comme  enseignant. Je donne
des  éléments d'appréciation  de la  candidature de  Julien à  ce titre.  Julien
Casarin, après ses classes préparatoires aux grandes écoles, a intégré l'ENSIIE,
via  le concours  Centrale-Supélec qui  était  alors le  concours d'entrée  pour
l'école. Il a  ensuite suivi un tronc commun (mêmes  enseignements et évaluation
sur les sites  d'Evry et de Strasbourg)  sur trois semestres.  Il  a validé sans
retard ces trois semestres.

Les semestres  4 et 5 à  Strasbourg sont systématiquement adossés  à des Masters
informatiques suivis en bi-cursus par  les élèves-ingénieurs. Les élèves suivent,
pendant ce cycle Master, 42 ECTS constitués uniquement d'enseignements techniques
essentiellement orientés recherche. Julien Casarin a suivi les enseignements du Master
\textit{Informatique   et   Sciences   de    l'Image}   porté   par   l'UFR   de
Mathématique-Informatique. Ces enseignements constituent une première approche
de la recherche dans ce domaine.

En  parallèle   à  l'ENSIIE,  Julien  Casarin   a  suivi  et  validé   deux  UEs
\textit{Projet R\&D} (6 ECTS,  150h) aux semestres 4 et 5  qui sont l'occasion de
mettre en  {\oe}uvre une  démarche de  type recherche, de  l'étude de  l'état de
l'art jusqu'à la réalisation d'un prototype. Il a travaillé au cours de ses deux
projets R\&D  : 
\begin{itemize}
\item sur  une partie d'un modeleur  3D pour la  fabrication d'objets
imprimés en  3D, selon un cahier  des charges de la  société \textit{Skimlab}, 
\item sur  une application mobile  permettant à  ses utilisateurs de  trouver un
  point de  rendez vous  optimal étant  donné un  graphe représentant  les lieux
  d'une ville et les moyens de transport.
\end{itemize}



Enfin, Julien Casarin a manifesté durant sa scolarité son goût pour la recherche
en rejoignant lors  de ses deux derniers  stages (dont le stage  de fin d'études
d'un semestre) l'entreprise Accelinn. Le coeur de métier de cette entreprise est
l'apport de  solutions innovantes  et son personnel  technique est  constitué en
majorité de docteurs dans différentes disciplines. Ayant fait partie de son jury
de soutenance,  je peux témoigner que  Julien a su s'intégrer  parfaitement dans
cette équipe et a pu contribuer utilement sur plusieurs missions.

 
Pour ces raisons, je suis convaincu que Julien Casarin est un excellent candidat
pour le cycle doctoral.



\end{letter}
\begin{flushright}
%\includegraphics[width=.2\textwidth]{signgenaud.jpg}
\end{flushright}
\end{document}
